\documentclass{article}
\usepackage{geometry}
\usepackage{concmath}
\usepackage[OT1]{fontenc}
\usepackage{graphicx}
\usepackage{amsfonts}
\usepackage{amsmath,mathtools}
\usepackage{calrsfs}
\usepackage{color}
\usepackage[dvipsnames]{xcolor}
\usepackage[many]{tcolorbox}
\usepackage{setspace}
\usepackage{tkz-tab}    
\usepackage{amssymb}          
\usepackage{multicol}  
\usepackage{esint}
\geometry{margin=0.5in}
\definecolor{green}{RGB}{52, 235, 168}
\definecolor{violet}{RGB}{165, 83, 237}
\definecolor{lemon}{RGB}{170, 237, 83}
\definecolor{red}{RGB}{240, 125, 84}
\definecolor{blue}{RGB}{105, 255, 242}

\begin{document}
\noindent \colorbox{green}{240.1}\\
\begin{align*}
    J &= \begin{pmatrix}
        1 & ... & 1 \\
        \vdots & & \vdots \\
        1 & ... & 1 
    \end{pmatrix} \in \mathcal{M}_n(\mathbb{R})\\
    \chi_J &= \begin{vmatrix}
        X-1 & -1 & ... & -1 \\
        -1 & X-1 & ... & \vdots \\ 
        \vdots & & \ddots & \vdots \\
        -1 & ... & ... & -1
    \end{vmatrix}\\
    &= \begin{vmatrix}
        X-n & -1 & ... & -1 \\
        X-n & X-1 & ... & \vdots\\
        \vdots & & \ddots & \vdots \\
        X-n & -1 & ... & X-1
    \end{vmatrix} C_1 \leftarrow C_1 + C_2 + ... + C_n \\
    &= (X-n) \begin{vmatrix}
        1 & -1 & ... & -1 \\
        \vdots & X-1 & & \vdots \\
        \vdots & -1 & \ddots & \vdots \\
        1 & -1 & ... & X-1
    \end{vmatrix}\\
    &= (X-n) \begin{vmatrix}
        1 & 0 & ... & 0 \\
        \vdots & X & & \vdots \\
        \vdots &  & \ddots & \vdots \\
        1 & ... & ... & X
    \end{vmatrix}\\
    &= (X-n)X^{n-1} \text{ car matrice triangulaire}\\
    Sp(J) &= \{0, n\}\\
    X = \begin{pmatrix}
        x_1\\\vdots\\x_n
    \end{pmatrix} \in E_0(J) &\Leftrightarrow JX = 0 \\
    &\Leftrightarrow \begin{cases}
        x_1 + ... + x_n = 0\\
        x_1 ... x_n = 0
    \end{cases}\\
    &\Leftrightarrow x_1 +...+x_n = 0\\
    &\Leftrightarrow x_n = -x_1 -...-x_{n-1}\\
    &\Leftrightarrow X = \begin{pmatrix}
        x_1\\\vdots\\-x_1-...-x_{n-1}
    \end{pmatrix} = x_1 \begin{pmatrix}
        1\\0\\\vdots\\0\\-1
    \end{pmatrix} +...+x_{n-1}\begin{pmatrix}
        0\\\vdots\\0\\1\\-1
    \end{pmatrix}
    \Leftrightarrow X \in Vect(U_1,...,U_{n-1})
\end{align*}
Donc \(E_0 (J) = Vect(U_1,...,U_{n-1})\)\\
\(E_0(J) = Ker(J-0I_n)\)\\

\noindent \colorbox{green}{240.3}\\
\begin{equation*}
    A = \begin{pmatrix}
        1&0\\0&2
    \end{pmatrix} \indent
    M = \begin{pmatrix}
        a&b\\c&d
    \end{pmatrix}
\end{equation*}
a/ \begin{align*}
    AM = \begin{pmatrix}
        a&b\\2c&2d
    \end{pmatrix} &\indent MA = \begin{pmatrix}
        a&2b\\c&2d
    \end{pmatrix}\\
    AM - MA &= \begin{pmatrix}
        a&-b\\c&0
    \end{pmatrix}
\end{align*}
b/ \begin{align*}
    \varphi: &\mathcal{M}_2(\mathbb{R}) \to \mathbb{M}_2(\mathbb{R})\\
    &M \mapsto AM - MA
\end{align*}
M2 \begin{align*}
    &M2 : AM - MA = \lambda M\\
    &\Leftrightarrow \begin{pmatrix}
        a&-b\\c&0
    \end{pmatrix} = \lambda \begin{pmatrix}
        a & b \\ c & d
    \end{pmatrix}\\
    &\Leftrightarrow \begin{cases}
        \lambda a = 0 \\ \lambda b = -b \\ \lambda c = c \\ \lambda d = 0
    \end{cases}
\end{align*}
1er cas \(\Leftrightarrow \lambda = 0\) \(
\begin{cases}
    b = 0 \\
    c = 0
\end{cases}\)\\
Donc \(E_0(\varphi) = \{\begin{pmatrix}
a & 0 \\ 0 & d
\end{pmatrix} | a,d \in \mathbb{R}\}\)\\
Si \(\lambda = 1\) a = b = d = 0 \\
\(E_1(\varphi) = \{\begin{pmatrix}
    0 & 0 \\ c & 0
\end{pmatrix} | c \in \mathbb{R}\}\)\\
Si \(\lambda = -1\) \\
\(E_{-1}(\varphi) = \{\begin{pmatrix}
    0 & b \\ 0 & 0
\end{pmatrix} | b \in \mathbb{R}\}\)\\

\noindent \colorbox{green}{240.2}\\
\begin{align*}
    \chi_c &= \begin{vmatrix}
        X & 0 & \dots & 0 & a_1\\
        -1 & X &  &   & \vdots\\
        0 & -1 & \ddots &  & \vdots\\
        \vdots & & \ddots & X & \vdots \\
        0 & \dots & 0 & -1 & -a_{n-1}
    \end{vmatrix} L_1 \leftarrow L_1 + XL_2 + X^2 L_3 + ... _ X^{n-1}L_n\\
    &= \begin{vmatrix}
        0 & ... & ...&... & 0 & P \\
        -1 & X & & & & a_1\\
        0 & -1 & X & & & \vdots\\
        & & \ddots & \ddots & & \vdots\\
        & & &\ddots  &  X & a_{n-2}\\
        0& & & 0& -1 & X + a_{n-1}
    \end{vmatrix}\\
    P &= a_0 + Xa_1 + X^2 a_2 + ... + X^{n-1}(X+a_{n-1})\\
    &\text{dev } L_1 = (-1)^{n+1}P(-1)^{n-1} = P
\end{align*}\\
b/ \begin{align*}
    V = V(\lambda_1,...,\lambda_n) &= \begin{pmatrix}
        1 & ... & 1\\
        \lambda_1 & ... & \lambda_n\\
        \vdots & & \vdots \\
        \lambda_1^{n-1}&...&\lambda_n^{n-1}
    \end{pmatrix}\\
    det(V(\lambda_1,...,\lambda_n)) &= \prod_{1 \leq i < j \leq n}^{}(\lambda_j - \lambda_i) \neq 0
\end{align*}
Donc \(V^{-1}\) existe\\
B = \(C_1...C_n\) avec \(V=(C_1|...|C_n)\)\\
V = Pass(\(\mathcal{C}\to B\))\begin{align*}
    u(C_k) = C^T C_k \\
    &= \begin{pmatrix}
        0 & 1 & & \\
        & \ddots & \ddots & & \\
        & & 0 & 1 \\
        -a_0&... & ...&-a_{n-1}
    \end{pmatrix}\begin{pmatrix}
        1 \\ \vdots \\ \vdots \\ \lambda_k^{n-1}
    \end{pmatrix}\\
    &= \begin{pmatrix}
        \lambda_k \\
        \lambda^2_k \\
        \vdots \\
        \lambda_k^{n-1}\\
        -a_0 - a_1\lambda_k - a_2\lambda_k^2-...-a_{n-1}\lambda_k^{n-1}
    \end{pmatrix}\\
    &= \begin{pmatrix}
        \lambda_k \\
        \lambda^2_k \\
        \vdots \\
        \lambda_k^{n-1}\\
        -P(\lambda_k) + \lambda_k^n
    \end{pmatrix}\\
    &= \lambda_k \begin{pmatrix}
        1 \\ \vdots \\ \lambda_k^{n-1}
    \end{pmatrix}
    &= \lambda_k C_k
\end{align*}
donc \begin{equation*}
    Mat(u,B) = \begin{pmatrix}
        \lambda_1 & & & \\
        & \ddots & & \\
        & & & \lambda_n
    \end{pmatrix}
\end{equation*}
\colorbox{green}{240.6}\\
\begin{align*}
    \chi_J = X^n -1
\end{align*}
\(\chi_J\) est scindé donc J diagonalisable\\
b/ u canoniquement associé à J dans \(\mathcal{B}(e_0,...,e_{n-1})\)\begin{equation*}
    u(e_j) = e_{j-1[n]}
\end{equation*} 
donc \(u^k(e_j) = e_{j-k[n]}\)\\
\(J^k=\) J où les colonnes sont décalées de k places modulo n a droite \\
\begin{align*}
    A &= A_0 I_n + a_1 J + ... + a_{n-1}J^{n-1}\\
    &= P(J) \text{ où } P = a_0 + a_1 X + ... + a_{n-1}J^{n-1}
\end{align*}
où \(J = Q D Q^{-1}\) avec \begin{align*}
    D &= \begin{pmatrix}
        1 & & & & \\
        & \omega & & & \\
        & & \omega^2 & & \\
        & & & \ddots & \\
        & & & & \omega^{n-1}
    \end{pmatrix}\\
    J&= Q \begin{pmatrix}
       P( 1) & & & & \\
        & P(\omega )& & & \\
        & & P(\omega^2 )& & \\
        & & & \ddots & \\
        & & & & P(\omega^{n-1})
    \end{pmatrix}Q^{-1}
\end{align*}
donc \begin{equation*}
    det A = \prod_{k=0}^{n-1}(\sum_{i=0}^{n-1}a_i \omega^{ki})
\end{equation*}
\noindent \colorbox{green}{240.4}\\
\begin{align*}
    A = \begin{pmatrix}
        6 & 0& 0\\
        0& 2& -2\\
        0 & 0 & 0
    \end{pmatrix}\\
    Sp(A) = \{0,2,6\}\\
    u(e_1) = 6e_1 
\end{align*}
Donc \(e_1 \in E_6(u)\) or \(dim E_6(u)\leq m_6 \leq 1\)\\
Donc \(E_6(u) = Vect(e_1)\)\\
de même \(E_2(u) = Vect(e_2)\)\\
\(u(e_2 + e_3) = 0\) et \(dimE_0(u) = 1\)\\
Donc \(E_0(u) = Vect(e_2 + e_3)\)\\
On pose \(D = \begin{pmatrix}
    6 & 0&0\\0&2&0\\0&0&0
\end{pmatrix} = Mat(u, B')\) et P = \(\begin{pmatrix}
    1&0&0\\0&1&1\\0&0&1
\end{pmatrix} = Pass(B \to B')\) par changement de base \fbox{\(A = PDP^{-1}\)}
c/ (E) \(X^2+X=A\) avec \(X \in \mathcal{M}_3(\mathbb{R})\)\\
\underline{Analyse}\\
Soit X tq \(X^2 + X = A\) \begin{align*}
    XA &= X(X^2 + X)\\
    &= X^3 + X^2\\
    &= (X^2 + X)X\\
    &= AX
\end{align*}
\(X^2 + X = A\) donc \(PY^2P^{-1} + PYP^{-1} = PDP^{-1}\)\\
donc \(Y^2 + Y = D\)
\begin{align*}
    &Y = \begin{pmatrix}
        \alpha & 0 & 0\\ 0& \beta  & 0 \\ 0& 0 & \gamma
    \end{pmatrix}\\
    &\text{ donc } \begin{pmatrix}
        \alpha^2 & 0 & 0\\ 0& \beta^2  & 0 \\ 0& 0 & \gamma^2
    \end{pmatrix} + \begin{pmatrix}
        \alpha & 0 & 0\\ 0& \beta  & 0 \\ 0& 0 & \gamma
    \end{pmatrix} = \begin{pmatrix}
        6 & 0 & 0\\ 0& 2  & 0 \\ 0& 0 & 0
    \end{pmatrix}\\
    &\text{ donc } \begin{cases}
        \alpha^2 + \alpha - 6 = 0\\
        \beta^2 + \beta - 2 = 0\\
        \gamma^2 + \gamma = 0
    \end{cases} \text{ cad } \begin{cases}
        (\alpha-2)(\alpha+3) = 0 \\
        (\beta-1)(\beta + 2) = 0\\
        \gamma(\gamma+1) = 0
    \end{cases}
\end{align*}
\colorbox{green}{240.31}\\
\begin{equation*}
    A \in \mathcal{M}_n(\mathcal{K}), \exists p \in \mathbb{N}, A^p = 0_n
\end{equation*}
Mq \(Sp(A) \subset \{0\}\)\\
Soit \(\lambda\) vp de A 
\begin{align*}
    \exists X \neq 0_{n,1}, AX &= \lambda X\\
    A^p X &= \lambda^p X \\
    \lambda^p X &= \begin{pmatrix}
        0\\ \vdots \\ 0
    \end{pmatrix} \text{ Donc } \lambda = 0\\
    \chi_A &= X^n\\
    \chi_A &= X^n - tr(A)X^{n-1}+...+(-1)^ndet(A)\\
    trA &= 0
\end{align*}
2/ \begin{align*}
    A &= \begin{pmatrix}
        3 & 1 & -1\\ 1& 1& 1\\ 2& 0 & 2
    \end{pmatrix}\\
    tr(A-\lambda Id_3) &= trA - \lambda tr(Id_3)\\
    &= 6 - 3\lambda\\
\end{align*}
nécessairement \(\lambda = 2\)\\
Mq \(B = \begin{pmatrix}
    1 & 1 & -1\\ 1& -1 & 1 \\ 2 & 0 & 0 
\end{pmatrix}\) est nilpotente 
\(B^3 = 0\)\\
\colorbox{green}{250.7}\\
a/ Mq \(P(A^T) = P(A)^T\)
\begin{align*}
    P &= \sum_{k-0}^{n}a_k X^k\\
    P(A^T) &= \sum_{k=0}^{n}a_k(A^T)^k\\
    P(A)^T &= (\sum_{k-0}^{n}a_k X^k)^T = \sum_{k=0}^{n}a_k (A^k)^T
\end{align*}
Mq : \((A^T)^k = (A^k)^T\)\\
Relation de récurrence : \((A^T)^{k+1} = (A^T)^k A^T = (AA^k)^T\)\\
b/ Contrexemple / vrai 



\end{document}
